\section*{T/F (sample)}
\textbf{(a) False.} Desktop operating systems generally use FCFS for CPU scheduling.\\
\textbf{(b) False.} Using operational analysis, given the mean think time and demands for devices in a client/server system, an exact value for throughput can be calculated.\\
\textbf{(c) False.} Round-robin routing yields lower mean response times than JSQ (Join the Shortest Queue).\\
\textbf{(d) True.} The Websphere availability study used a mixture of fault tree and CTMC (Continuous-Time Markov Chain) models.\\
\textbf{(e) True.} PS (Processor Sharing) and LCFS (Last-Come, First-Served) have the same mean response time in an $M/G/1$ queue.\\
\textbf{(f) False.} In a multiserver system, one needs approximately $2\lambda/\mu$ servers to ensure that less than 20 percent of arrivals have to wait.\\
\textbf{(g) False.} Little's Law only applies for $M/M/1$ queues.\\
\textbf{(h) False.} The bottleneck in a network is the node that is visited most frequently.\\
\textbf{(i) False.} Using MVA (Mean Value Analysis), one can compute the likelihood that a node is idle.\\
\textbf{(j) True.} For a random variable $X$, $E[X^2]$ is always greater than or equal to $\Var(X)$.
